\documentclass[../../main.tex]{subfiles}% 注意这里的文件路径不能用 ./main.tex ,否则用latexmk编译子文件会报错
\graphicspath{{\subfix{./image/}}{\subfix{../../image/}}} % 指定图片引用路径,后续可以直接使用图片文件名
% 如果图片存放在上级目录,则使用 ../../image/ 作为备用路径

% 例如:
% \begin{figure}[H]
% \centering
% \includegraphics[width=0.3\textwidth]{图.png}
% \caption{}
% \label{figure:图}
% \end{figure}
% 注意:上述\label{}一定要放在\caption{}之后,否则引用图片序号会只会显示??.

\begin{document}
% SOURCE p=100; printed=85; object=ch03-signed-measures-title; region=middle
\section{符号测度}\label{section:ch03-signed-measures}
% SOURCE p=100; printed=85; object=ch03-signed-measure; region=lower
\begin{definition}[符号测度]\label{definition:ch03-signed-measure}
设 $(X,\mathscr{M})$ 为可测空间。$(X,\mathscr{M})$ 上的\textbf{符号测度}是函数 $\nu:\mathscr{M}\to[-\infty,\infty]$，满足：
\begin{enumerate}
\item $\nu(\varnothing)=0$。
\item $\nu$ 在 $\infty$ 和 $-\infty$ 中至多取其中一个值。
\item 若 $\{E_j\}$ 为 $\mathscr{M}$ 中两两不交的集合序列，则
\begin{align*}
\nu\left(\bigcup_1^\infty E_j\right)=\sum_1^\infty\nu(E_j),
\end{align*}
并且当 $\nu(\bigcup_1^\infty E_j)$ 有限时，右端级数绝对收敛。
\end{enumerate}
\end{definition}
% SOURCE p=100; printed=85; object=ch03-positive-measure; region=bottom
\begin{definition}[正测度]\label{definition:ch03-positive-measure}
每个测度都是符号测度；在需要强调两者区别时，称测度为\textbf{正测度}。
\end{definition}
% SOURCE p=101; printed=86; object=ch03-signed-measure-examples; region=top
这里给出两个简单例子。
\begin{example}\label{example:ch03-difference-measures}
若 $\mu_1,\mu_2$ 为 $\mathscr{M}$ 上的正测度，且其中至少一个有限，则 $\nu=\mu_1-\mu_2$ 为符号测度。
\end{example}
\begin{example}\label{example:ch03-signed-density}
若 $\mu$ 为 $\mathscr{M}$ 上的正测度，$f:X\to[-\infty,\infty]$ 为可测函数，且 $\int f^+\,\mathrm{d}\mu$ 与 $\int f^-\,\mathrm{d}\mu$ 中至少一个有限，则 $\nu(E)=\int_E f\,\mathrm{d}\mu$ 为符号测度。
\begin{definition*}[广义可积]
满足上述条件的 $f$ 称为\textbf{广义 $\mu$-可积}的。
\end{definition*}
\end{example}
实际上，我们很快会看到，每个符号测度都可以表示为这两种形式中的任一种。
% SOURCE p=101; printed=86; object=ch03-signed-continuity; number=3.1; region=upper
\begin{proposition}\label{proposition:ch03-signed-continuity}
设 $\nu$ 为 $(X,\mathscr{M})$ 上的符号测度。
\begin{enumerate}
\item\label{item:ch03-signed-continuity-below} 若 $\{E_j\}$ 为 $\mathscr{M}$ 中的递增序列，则 $\nu(\bigcup_1^\infty E_j)=\lim_{j\to\infty}\nu(E_j)$。
\item\label{item:ch03-signed-continuity-above} 若 $\{E_j\}$ 为 $\mathscr{M}$ 中的递减序列，且 $\nu(E_1)$ 有限，则 $\nu(\bigcap_1^\infty E_j)=\lim_{j\to\infty}\nu(E_j)$。
\end{enumerate}
\end{proposition}
\begin{proof}
证明与正测度的情形（\cref{theorem:ch01-measure-properties}）基本相同，留给读者完成（\cref{exercise:ch03-signed-continuity-proof}）。
\end{proof}
% SOURCE p=101; printed=86; object=ch03-sign-sets; region=middle
\begin{definition}[正集、负集与零集]\label{definition:ch03-sign-sets}
设 $\nu$ 为 $(X,\mathscr{M})$ 上的符号测度。若对每个满足 $F\subseteq E$ 的 $F\in\mathscr{M}$，均有 $\nu(F)\ge0$、$\nu(F)\le0$ 或 $\nu(F)=0$，则分别称 $E\in\mathscr{M}$ 为关于 $\nu$ 的\textbf{正集}、\textbf{负集}或\textbf{零集}。
\end{definition}
% SOURCE p=101; printed=86; object=ch03-sign-sets-density; region=middle
\begin{example}\label{example:ch03-sign-sets-density}
若 $\nu(E)=\int_E f\,\mathrm{d}\mu$，则 $E$ 为正集、负集或零集，当且仅当在 $E$ 上分别有 $f\ge0$、$f\le0$ 或 $f=0$，$\mu$-几乎处处成立。
\end{example}
% SOURCE p=101; printed=86; object=ch03-positive-union; number=3.2; region=middle
\begin{lemma}\label{lemma:ch03-positive-union}
正集的任意可测子集都是正集，可数个正集的并也是正集。
\end{lemma}
\begin{proof}
第一个断言显然成立。若 $P_1,P_2,\ldots$ 为正集，令 $Q_n=P_n\setminus\bigcup_1^{n-1}P_j$。则 $Q_n\subseteq P_n$，故 $Q_n$ 为正集。若 $E\in\mathscr{M}$ 且 $E\subseteq\bigcup_1^\infty P_j$，则
\begin{align*}
\nu(E)=\sum_1^\infty\nu(E\cap Q_j)\ge0.
\end{align*}
\end{proof}
% SOURCE p=101; printed=86; object=ch03-hahn-decomposition; number=3.3; region=lower
\begin{theorem}[Hahn 分解定理]\label{theorem:ch03-hahn-decomposition}
若 $\nu$ 为 $(X,\mathscr{M})$ 上的符号测度，则存在一个正集 $P$ 和一个负集 $N$，使 $P\cup N=X$ 且 $P\cap N=\varnothing$。若 $P',N'$ 为另一对这样的集合，则 $P\mathbin{\triangle}P'$（$=N\mathbin{\triangle}N'$）为关于 $\nu$ 的零集。
\end{theorem}
\begin{proof}
假设 $\nu$ 从不取值 $\infty$；否则用 $-\nu$ 代替 $\nu$。令 $m=\sup\{\nu(E):E\text{ 为正集}\}$，取正集 $P_j$，使 $\nu(P_j)\to m$。由\cref{lemma:ch03-positive-union,proposition:ch03-signed-continuity}，$P=\bigcup_1^\infty P_j$ 为正集，且 $\nu(P)=m<\infty$。令 $N=X\setminus P$。我们断言 $N$ 为负集。

为此，假设 $N$ 不是负集，导出矛盾。首先，$N$ 不包含任何不是零集的正集 $E$，否则 $E\cup P$ 为正集，且 $\nu(E\cup P)=\nu(E)+\nu(P)>m$，这是不可能的。其次，若 $A\subseteq N$ 且 $\nu(A)>0$，则存在 $B\subseteq A$，使 $\nu(B)>\nu(A)$：因为 $A$ 不是正集，所以存在 $C\subseteq A$，使 $\nu(C)<0$；取 $B=A\setminus C$，便有 $\nu(B)=\nu(A)-\nu(C)>\nu(A)$。于是，若 $N$ 不是负集，可以按如下方式选取 $A_j\subseteq N$ 和正整数 $n_j$：$n_1$ 是使得存在 $B\subseteq N$ 满足 $\nu(B)>n_1^{-1}$ 的最小整数，$A_1$ 就取为这样的一个 $B$；
% SOURCE p=102; printed=87; object=ch03-hahn-decomposition-proof-cont; region=top
依次令 $n_j$ 为使得存在 $B\subseteq A_{j-1}$ 满足 $\nu(B)>\nu(A_{j-1})+n_j^{-1}$ 的最小整数，$A_j$ 就取为这样的一个 $B$。令 $A=\bigcap_1^\infty A_j$。则
\begin{align*}
\infty>\nu(A)=\lim_{j\to\infty}\nu(A_j)>\sum_1^\infty n_j^{-1},
\end{align*}
所以 $n_j\to\infty$。但是，再次应用上面的论证，就存在 $B\subseteq A$ 和某个 $n\in\mathbb{N}$，使 $\nu(B)>\nu(A)+n^{-1}$。当 $j$ 充分大时，$n<n_j$，且 $B\subseteq A_{j-1}$，这与 $n_j$ 的选取矛盾。因此 $N$ 为负集。

最后，若 $P',N'$ 为另一对这样的集合，则 $P\setminus P'\subseteq P$，同时 $P\setminus P'\subseteq N'$，所以它既是正集又是负集，因而是零集。$P'\setminus P$ 的情形相同，结论即得。
\end{proof}
% SOURCE p=102; printed=87; object=ch03-hahn-decomposition-definition; region=upper
\begin{definition}[Hahn 分解]\label{definition:ch03-hahn-decomposition}
$X$ 分为不交的正集 $P$ 和负集 $N$ 的分解 $X=P\cup N$，称为关于 $\nu$ 的一个\textbf{Hahn 分解}。
\end{definition}
\begin{remark}
Hahn 分解通常不唯一，因为零集可以从 $P$ 移到 $N$，或反过来。不过，它给出一种将 $\nu$ 表示为两个正测度之差的典范方式。
\end{remark}
% SOURCE p=102; printed=87; object=ch03-mutually-singular; region=upper-middle
为叙述这一结果，我们需要一个新概念。
\begin{definition}[相互奇异]\label{definition:ch03-mutually-singular}
$(X,\mathscr{M})$ 上的两个符号测度 $\mu,\nu$ 称为\textbf{相互奇异}，若存在 $E,F\in\mathscr{M}$，使 $E\cap F=\varnothing$、$E\cup F=X$，且 $E$ 为 $\mu$-零集，$F$ 为 $\nu$-零集。也称 $\nu$ 关于 $\mu$ \textbf{奇异}，或反过来说；用垂直符号记作 $\mu\perp\nu$。
\end{definition}
\begin{remark}
粗略地说，相互奇异的测度“存在于”彼此不交的集合上。
\end{remark}
% SOURCE p=102; printed=87; object=ch03-jordan-decomposition; number=3.4; region=middle
\begin{theorem}[Jordan 分解定理]\label{theorem:ch03-jordan-decomposition}
若 $\nu$ 为符号测度，则存在唯一的正测度 $\nu^+,\nu^-$，满足 $\nu=\nu^+-\nu^-$ 且 $\nu^+\perp\nu^-$。
\end{theorem}
\begin{proof}
设 $X=P\cup N$ 为 $\nu$ 的 Hahn 分解，定义 $\nu^+(E)=\nu(E\cap P)$、$\nu^-(E)=-\nu(E\cap N)$。显然，这两个测度满足所需条件。

若还有 $\nu=\mu^+-\mu^-$，其中 $\mu^+\perp\mu^-$，取彼此不交且并为 $X$ 的集合 $E,F$，使 $\mu^+(F)=\mu^-(E)=0$。则 $X=E\cup F$ 为另一个 Hahn 分解，所以 $P\mathbin{\triangle}E$ 为 $\nu$-零集。因此，对 $A\in\mathscr{M}$，
\begin{align*}
\mu^+(A)=\mu^+(A\cap E)=\nu(A\cap E)=\nu(A\cap P)=\nu^+(A).
\end{align*}
同样，$\mu^-=\nu^-$。
\end{proof}
% SOURCE p=102; printed=87; object=ch03-jordan-variations; region=lower
\begin{definition}[正变差、负变差与 Jordan 分解]\label{definition:ch03-jordan-variations}
测度 $\nu^+,\nu^-$ 分别称为 $\nu$ 的\textbf{正变差}和\textbf{负变差}；表达式 $\nu=\nu^+-\nu^-$ 称为 $\nu$ 的\textbf{Jordan 分解}。
\end{definition}
这个名称源于它与有界变差实函数分解为两个递增函数之差的类比，后者将在\cref{section:ch03-bounded-variation}中讨论。
% SOURCE p=102; printed=87; object=ch03-total-variation; region=lower
\begin{definition}[全变差]\label{definition:ch03-total-variation}
$\nu$ 的\textbf{全变差}是测度 $|\nu|\triangleq\nu^++\nu^-$。
\end{definition}
% SOURCE p=102; printed=87; object=ch03-variation-null-singular; region=lower
\begin{proposition}\label{proposition:ch03-variation-null-singular}
$E$ 为 $\nu$-零集，当且仅当 $|\nu|(E)=0$。并且，$\nu\perp\mu$ 当且仅当 $|\nu|\perp\mu$，也当且仅当 $\nu^+\perp\mu$ 且 $\nu^-\perp\mu$。
\end{proposition}
\begin{proof}
这些等价关系容易验证（\cref{exercise:ch03-variation-null-singular}）。
\end{proof}
% SOURCE p=102; printed=87; object=ch03-signed-boundedness; region=bottom
\begin{proposition}\label{proposition:ch03-signed-boundedness}
若 $\nu$ 不取值 $\infty$，则 $\nu^+$ 有限，且 $\nu$ 有上界；若 $\nu$ 不取值 $-\infty$，则 $\nu^-$ 有限，且 $\nu$ 有下界。特别地，若 $\nu$ 为实值符号测度，则它有界。
\end{proposition}
\begin{proof}
若 $\nu$ 不取值 $\infty$，则 $\nu^+(X)=\nu(P)<\infty$，从而 $\nu(E)\le\nu^+(X)$。不取值 $-\infty$ 的情形同理。
\end{proof}
% SOURCE p=103; printed=88; object=ch03-total-variation-density; region=top
\begin{proposition}\label{proposition:ch03-total-variation-density}
每个符号测度都可以表示为 $\nu(E)=\int_E f\,\mathrm{d}\mu$，其中 $\mu=|\nu|$，$f=\chi_P-\chi_N$，而 $X=P\cup N$ 是 $\nu$ 的一个 Hahn 分解。
\end{proposition}
% SOURCE p=103; printed=88; object=ch03-signed-integration; region=top
\begin{definition}[关于符号测度的积分]\label{definition:ch03-signed-integration}
定义
\begin{align*}
L^1(\nu)&\triangleq L^1(\nu^+)\cap L^1(\nu^-),\\
\int f\,\mathrm{d}\nu&\triangleq\int f\,\mathrm{d}\nu^+-\int f\,\mathrm{d}\nu^-\qquad(f\in L^1(\nu)).
\end{align*}
\end{definition}
% SOURCE p=103; printed=88; object=ch03-finite-signed-measure; region=upper
\begin{definition}[有限与 \texorpdfstring{$\sigma$}{σ}-有限符号测度]\label{definition:ch03-finite-signed-measure}
若 $|\nu|$ 有限或 $\sigma$-有限，则分别称 $\nu$ 为\textbf{有限}或\textbf{$\sigma$-有限}符号测度。
\end{definition}
% SOURCE p=103; printed=88; object=ch03-signed-exercises-title; region=upper
\subsection*{习题}
% SOURCE p=103; printed=88; object=ch03-ex01; number=1; region=upper
\begin{exercise}\label{exercise:ch03-signed-continuity-proof}
证明\cref{proposition:ch03-signed-continuity}。
\end{exercise}
% SOURCE p=103; printed=88; object=ch03-ex02; number=2; region=upper
\begin{exercise}\label{exercise:ch03-variation-null-singular}
若 $\nu$ 为符号测度，则 $E$ 为 $\nu$-零集当且仅当 $|\nu|(E)=0$。此外，若 $\nu,\mu$ 为符号测度，则 $\nu\perp\mu$ 当且仅当 $|\nu|\perp\mu$，也当且仅当 $\nu^+\perp\mu$ 且 $\nu^-\perp\mu$。
\end{exercise}
% SOURCE p=103; printed=88; object=ch03-ex03; number=3; region=upper-middle
\begin{exercise}\label{exercise:ch03-variation-integral}
设 $\nu$ 为 $(X,\mathscr{M})$ 上的符号测度。
\begin{enumerate}
\item $L^1(\nu)=L^1(|\nu|)$。
\item 若 $f\in L^1(\nu)$，则 $|\int f\,\mathrm{d}\nu|\le\int|f|\,\mathrm{d}|\nu|$。
\item 若 $E\in\mathscr{M}$，则 $|\nu|(E)=\sup\{|\int_E f\,\mathrm{d}\nu|:|f|\le1\}$。
\end{enumerate}
\end{exercise}
% SOURCE p=103; printed=88; object=ch03-ex04; number=4; region=middle
\begin{exercise}\label{exercise:ch03-jordan-minimal}
若 $\nu$ 为符号测度，$\lambda,\mu$ 为正测度，且 $\nu=\lambda-\mu$，则 $\lambda\ge\nu^+$ 且 $\mu\ge\nu^-$。
\end{exercise}
% SOURCE p=103; printed=88; object=ch03-ex05; number=5; region=middle
\begin{exercise}\label{exercise:ch03-variation-triangle}
若 $\nu_1,\nu_2$ 为符号测度，且两者都不取值 $\infty$，或两者都不取值 $-\infty$，则 $|\nu_1+\nu_2|\le|\nu_1|+|\nu_2|$。（使用\cref{exercise:ch03-jordan-minimal}。）
\end{exercise}
% SOURCE p=103; printed=88; object=ch03-ex06; number=6; region=middle
\begin{exercise}\label{exercise:ch03-density-variations}
设 $\mu$ 为正测度，$f$ 为广义 $\mu$-可积函数，$\nu(E)=\int_E f\,\mathrm{d}\mu$。用 $f$ 和 $\mu$ 描述 $\nu$ 的 Hahn 分解以及正变差、负变差和全变差。
\end{exercise}
% SOURCE p=103; printed=88; object=ch03-ex07; number=7; region=middle-lower
\begin{exercise}\label{exercise:ch03-variation-set-suprema}
设 $\nu$ 为 $(X,\mathscr{M})$ 上的符号测度，$E\in\mathscr{M}$。
\begin{enumerate}
\item $\nu^+(E)=\sup\{\nu(F):F\in\mathscr{M},F\subseteq E\}$，$\nu^-(E)=-\inf\{\nu(F):F\in\mathscr{M},F\subseteq E\}$。
\item $|\nu|(E)=\sup\{\sum_1^n|\nu(E_j)|:n\in\mathbb{N},E_j\text{ 两两不交},\bigcup_1^nE_j=E\}$。
\end{enumerate}
\end{exercise}
\end{document}
